\documentclass[runningheads]{llncs}

\usepackage{amsmath,amssymb,mathtools}
\usepackage{microtype}
\usepackage{booktabs}
\usepackage{xcolor}
\usepackage{hyperref}

\hypersetup{
  colorlinks=true,
  linkcolor=blue,
  citecolor=blue,
  urlcolor=blue
}

\usepackage{listings}
\lstdefinelanguage{lean}{
  keywords={theorem,lemma,def,inductive,structure,namespace,end,section,variable,
    by,fun,match,with,have,show,let,in,where,do,if,then,else,intro,cases,simp,
    exact,apply,omega,sorry,import,open,deriving},
  sensitive=true,
  morecomment=[l]{--},
  morecomment=[s]{/-}{-/},
  morestring=[b]"
}
\lstset{
  language=lean,
  basicstyle=\ttfamily\small,
  columns=fullflexible,
  breaklines=true,
  frame=single,
  rulecolor=\color{black!20},
  keywordstyle=\bfseries,
  commentstyle=\itshape\color{black!50},
  numbers=left,
  numberstyle=\tiny\color{black!40},
  numbersep=5pt,
  literate={∀}{{\ensuremath{\forall}}}1 {¬}{{\ensuremath{\lnot}}}1
           {→}{{\ensuremath{\rightarrow}}}1 {←}{{\ensuremath{\leftarrow}}}1
           {≥}{{\ensuremath{\geq}}}1 {≤}{{\ensuremath{\leq}}}1
}

\begin{document}

\title{A Proof Governance Kernel\\with Machine-Checked Safety Invariants}
\titlerunning{Proof Governance Kernel}

\author{Croft Adams\inst{1}}
\authorrunning{C.\ Adams}
\institute{Independent Research \\ \email{theadamsfamily1981@gmail.com}}

\maketitle

\begin{abstract}
Long-horizon mathematical proof campaigns routinely involve claims at
different evidence levels: some unconditionally proved, others conditional on
open conjectures, others supported only by heuristic computation.
Misclassifying conditional or heuristic claims as unconditional is a
common workflow failure mode in both human and AI-assisted mathematics.

We present a \emph{proof governance kernel}: a small, typed state machine that
tracks domain lifecycle (active, dropped, locked, reactivated), certificate
strength (proof, reduction, computation, conditional, heuristic), and lock
permanence (proven, conditional, heuristic). We formalize the kernel in Lean~4
and machine-check five safety invariants: (1)~proven locks are permanent,
(2)~soft certificates cannot produce proven locks, (3)~hard certificates are
strictly stronger than soft ones, (4)~every proven lock requires prior
evidence, and (5)~no lifecycle phase can be skipped.

We demonstrate the kernel on two case studies: tracking conditional versus
unconditional status in a number-theoretic proof campaign (the generalized
Fermat equation), and preventing circular dependencies in proof DAGs.
Unlike build-level or session-level governance, our kernel operates at
claim-level granularity, governing epistemic status across heterogeneous
evidence sources. All Lean~4 sources compile without errors and are
provided as a reproducibility artifact.

\keywords{proof engineering \and Lean~4 \and formal verification \and
governance \and certificate systems \and soundness}
\end{abstract}

%% ============================================================
\section{Introduction}
\label{sec:intro}

Mathematical research, whether conducted by humans or AI systems, proceeds
through claims of varying confidence. A conjecture supported by numerical
evidence differs fundamentally from a peer-reviewed theorem, yet both may
coexist in a single research programme. When proofs span months or years and
involve multiple contributors---increasingly including machine
learning systems whose outputs may be plausible but
unsound~\cite{FarrugiaEtAl2024}---the risk of \emph{unsound claim
promotion} becomes acute: a conditional result is inadvertently
treated as unconditional, or a heuristic observation is locked
as settled.

We propose a lightweight \emph{proof governance kernel} that addresses this
risk through three mechanisms:

\begin{enumerate}
\item \textbf{Typed certificates.} Every state transition carries a certificate
  whose type (proof, reduction, computation, conditional, or heuristic) is
  machine-checked. Certificate types are ordered by strength and partitioned
  into \emph{hard} (proof, reduction, computation) and \emph{soft}
  (conditional, heuristic).

\item \textbf{Typed locks.} When a domain is resolved, it receives a lock whose
  type (proven, conditional, heuristic) determines whether it can be reopened.
  The lock type must be compatible with the certificate type: only hard
  certificates can produce proven locks.

\item \textbf{Lifecycle enforcement.} A domain must pass through mandatory
  intermediate states (active $\to$ dropped $\to$ locked), preventing shortcuts
  that bypass evidence requirements.
\end{enumerate}

\noindent
We formalize this kernel in Lean~4 (v4.14.0) and prove five safety theorems
entirely within Lean's kernel, requiring no axioms beyond the Lean~4 type
theory. The formalization comprises approximately 350 lines of Lean across
three files.

\paragraph{Contributions.}
\begin{itemize}
\item A formal model of proof state governance with typed certificates
  and locks (\S\ref{sec:model}).
\item Five machine-checked safety invariants (\S\ref{sec:theorems}).
\item Two case studies demonstrating practical application
  (\S\ref{sec:case1}, \S\ref{sec:case2}).
\item A reproducibility artifact with build instructions (\S\ref{sec:artifact}).
\end{itemize}

%% ============================================================
\section{The Kernel Model}
\label{sec:model}

\subsection{Certificate Types}

We define five certificate types, ordered by decreasing strength:

\begin{definition}[Certificate Type]
\label{def:certtype}
A \emph{certificate type} is one of five values, ordered by
decreasing strength: \texttt{proof}~(5), \texttt{reduction}~(4),
\texttt{computation}~(3), \texttt{conditional}~(2),
\texttt{heuristic}~(1). A certificate type is \emph{hard} if it
is \texttt{proof}, \texttt{reduction}, or \texttt{computation};
otherwise it is \emph{soft}.
\end{definition}

\noindent
In Lean~4:

\begin{lstlisting}
inductive CertificateType where
  | proof | reduction | computation | conditional | heuristic

def CertificateType.strength : CertificateType → Nat
  | .proof => 5 | .reduction => 4 | .computation => 3
  | .conditional => 2 | .heuristic => 1

def CertificateType.isHard : CertificateType → Bool
  | .proof => true | .reduction => true | .computation => true
  | .conditional => false | .heuristic => false
\end{lstlisting}

\subsection{Lock Types}

\begin{definition}[Lock Type]
\label{def:locktype}
A \emph{lock type} is one of:
\begin{itemize}
\item \texttt{proven}: permanent (cannot reopen),
\item \texttt{conditional}: reopenable upon condition resolution,
\item \texttt{heuristic}: reopenable with new tools.
\end{itemize}
\end{definition}

\begin{lstlisting}
def LockType.canReopen : LockType → Bool
  | .proven => false
  | .conditional => true
  | .heuristic => true
\end{lstlisting}

\subsection{Domain States and Transitions}

\begin{definition}[Domain State]
\label{def:state}
A domain occupies one of four states: \texttt{active} (under investigation),
\texttt{dropped}$(ct)$ (reduced with certificate type $ct$),
\texttt{locked}$(lt)$ (resolved with lock type $lt$), or \texttt{reactivated}
(reopened from a non-proven lock).
\end{definition}

\begin{definition}[Valid Transition]
\label{def:transition}
The permitted transitions are defined inductively:
\begin{align}
  &\texttt{active} \to \texttt{dropped}(ct) &&\text{for any } ct \label{eq:t1}\\
  &\texttt{dropped}(ct) \to \texttt{locked}(lt) &&\text{for any } lt \label{eq:t2}\\
  &\texttt{locked}(lt) \to \texttt{reactivated} &&\text{iff } lt.\texttt{canReopen} = \texttt{true} \label{eq:t3}\\
  &\texttt{reactivated} \to \texttt{active} && \label{eq:t4}
\end{align}
\end{definition}

\subsection{Lock Permission}

\begin{definition}[Permitted Lock]
\label{def:permitted}
A certificate type $ct$ \emph{permits} a lock type $lt$ according to:
\begin{itemize}
\item $\texttt{permittedLock}(ct, \texttt{proven})$ iff $ct.\texttt{isHard} = \texttt{true}$;
\item $\texttt{permittedLock}(ct, \texttt{conditional})$ iff $ct = \texttt{conditional}$ or $ct.\texttt{isHard} = \texttt{true}$;
\item $\texttt{permittedLock}(ct, \texttt{heuristic})$ always.
\end{itemize}
\end{definition}

%% ============================================================
\section{Safety Theorems}
\label{sec:theorems}

All governance kernel theorems below are machine-checked in Lean~4
with no \texttt{sorry} and no additional axioms. (The case study in
\S\ref{sec:case1} includes one intentional \texttt{sorry} marking
conditional analytic content; it is not used by any kernel theorem.)

\begin{theorem}[Proven Lock Permanence]
\label{thm:permanent}
No valid transition exits the proven-locked state:
\[
\forall\, s.\; \neg\, \texttt{ValidTransition}(\texttt{locked}(\texttt{proven}),\, s).
\]
\end{theorem}

\begin{proof}
By case analysis on \texttt{ValidTransition}. The only constructor producing a
transition from $\texttt{locked}(lt)$ is \texttt{lockToReactivate}, which
requires $lt.\texttt{canReopen} = \texttt{true}$. Since
$\texttt{proven}.\texttt{canReopen} = \texttt{false}$, no constructor applies.
\end{proof}

\begin{theorem}[No Premature Proven Lock]
\label{thm:nopremature}
Soft certificates cannot produce proven locks:
\begin{align*}
  &\neg\, \texttt{permittedLock}(\texttt{heuristic},\, \texttt{proven}) \\
  &\neg\, \texttt{permittedLock}(\texttt{conditional},\, \texttt{proven})
\end{align*}
\end{theorem}

\begin{proof}
By unfolding \texttt{permittedLock} and \texttt{isHard}:
both \texttt{heuristic.isHard} and \texttt{conditional.isHard}
reduce to \texttt{false}, contradicting the requirement for a
proven lock.
\end{proof}

\begin{theorem}[Strength Gap]
\label{thm:gap}
Every hard certificate type has strictly greater strength than every soft
certificate type:
\begin{multline*}
\forall\, ct_h\, ct_s.\;
  ct_h.\texttt{isHard} = \texttt{true} \;\land\;
  ct_s.\texttt{isHard} = \texttt{false} \\
  \Longrightarrow\;
  ct_h.\texttt{strength} > ct_s.\texttt{strength}.
\end{multline*}
\end{theorem}

\begin{proof}
Hard types have strength $\geq 3$ and soft types have strength $\leq 2$,
established by exhaustive case analysis on the five constructors.
\end{proof}

\begin{theorem}[Evidence Requirement]
\label{thm:evidence}
Any transition chain from \texttt{active} to $\texttt{locked}(\texttt{proven})$
must include an intermediate \texttt{dropped} state:
\begin{multline*}
\texttt{CertChain}(\texttt{active},\, \texttt{locked}(\texttt{proven})) \implies \\
  \exists\, ct.\; \texttt{Nonempty}(\texttt{CertChain}(\texttt{active},\, \texttt{dropped}(ct))).
\end{multline*}
\end{theorem}

\begin{proof}
By case analysis on the chain. A single-step chain from
\texttt{active} to \texttt{locked\allowbreak(proven)} is impossible
(no matching constructor). A multi-step chain must begin
with \texttt{activeToDrop}$(ct)$, giving the witness.
\end{proof}

\begin{theorem}[No State Skipping]
\label{thm:noskip}
The \texttt{active} state cannot transition directly to any locked state:
$$\forall\, lt.\; \neg\, \texttt{ValidTransition}(\texttt{active},\, \texttt{locked}(lt)).$$
\end{theorem}

\begin{proof}
By case analysis: no constructor of \texttt{ValidTransition} has source
\texttt{active} and target $\texttt{locked}(lt)$.
\end{proof}

%% ============================================================
\section{Case Study I: Conditional Certificates in Number Theory}
\label{sec:case1}

We instantiate the governance kernel on the generalized Fermat equation
$a^p + b^q = c^r$, specifically tracking the distinction between conditional
and unconditional results.

\paragraph{Background.}
Define $\lambda = 1/p + 1/q + 1/r$. When $\lambda < 1$, the equation has at
most finitely many coprime solutions by the Darmon--Granville theorem
\cite{DarmonGranville1995}, which invokes Faltings' theorem
\cite{Faltings1983}. A standard application of the ABC conjecture
\cite{Masser1985,Oesterle1988} yields a shorter finiteness argument, but this
is \emph{conditional} on the unproved ABC conjecture.

\paragraph{Certificate Assignment.}
In our kernel, the Darmon--Granville result receives a \texttt{proof}
certificate (unconditional, peer-reviewed). The ABC-based argument receives a
\texttt{conditional} certificate with an explicit blocking seam: ``ABC
conjecture unproved.''

\paragraph{Governance Enforcement.}
We verify in Lean~4:
\begin{itemize}
\item The ABC certificate is not hard:
  $\texttt{abcBealCertificate}.\texttt{isHard} = \texttt{false}$.
\item It cannot support a proven lock:\\
  $\neg\, \texttt{permittedLock}(\texttt{conditional},\, \texttt{proven})$.
\item The Darmon--Granville certificate \emph{can} support a proven lock:\\
  $\texttt{permittedLock}(\texttt{proof},\, \texttt{proven})$.
\end{itemize}

\paragraph{Signature Verification.}
We additionally verify regime classifications for specific signatures using
Lean's \texttt{native\_decide} tactic:
\begin{itemize}
\item $(2,2,2)$: spherical ($\lambda = 3/2 > 1$).
\item $(3,3,3)$, $(2,3,6)$, $(2,4,4)$: euclidean ($\lambda = 1$).
\item $(3,5,7)$: hyperbolic ($\lambda = 71/105 < 1$).
\end{itemize}

%% ============================================================
\section{Case Study II: Circular Dependency Prevention}
\label{sec:case2}

A second failure mode in proof campaigns is \emph{circular dependency}: claim
$A$ depends on claim $B$, which depends on $A$. Our kernel prevents this
structurally.

\paragraph{Structural Argument.}
Define a rank function: $\texttt{active} \mapsto 0$,
$\texttt{dropped}(\_) \mapsto 1$, $\texttt{locked}(\_) \mapsto 2$,
$\texttt{reactivated} \mapsto 0$.
By Theorem~\ref{thm:noskip}, forward transitions (excluding
reactivation) strictly increase rank. Since the rank space is
finite, any forward-only chain terminates. Reactivation resets
rank to zero but requires the lock to be non-proven
(Theorem~\ref{thm:permanent}).

\paragraph{Mechanized Impossibility.}
We verify in Lean~4 that the proven-locked state is a dead end:

\begin{lstlisting}[mathescape=true]
-- No transition leaves proven-locked state
theorem proven_is_terminal :
    $\forall$ (s : DomainState),
    $\lnot$ ValidTransition (.locked .proven) s := by
  intro s h; cases h with
  | lockToReactivate lt hcan =>
    simp [LockType.canReopen] at hcan

-- Cannot skip active directly to locked
theorem no_skip_active_to_locked :
    $\forall$ (lt : LockType),
    $\lnot$ ValidTransition .active (.locked lt) := by
  intro lt h; cases h
\end{lstlisting}

\begin{corollary}[No Proven Lock Cycle]
\label{cor:nocycle}
A domain that reaches $\texttt{locked}(\texttt{proven})$ cannot
return to \texttt{active} via any sequence of valid transitions.
\end{corollary}

\begin{proof}
Immediate from \texttt{proven\_is\_terminal}: no constructor of
\texttt{ValidTransition} matches the source
$\texttt{locked}(\texttt{proven})$.
\end{proof}

\paragraph{Conditional and Heuristic Cycles.}
Conditional and heuristic locks \emph{do} permit cycles
(active $\to$ dropped $\to$ locked $\to$ reactivated $\to$
active), which is intentional: new tools or resolved conditions
should allow revisiting a domain. The kernel guarantees that
such cycles never produce proven locks from soft evidence
(Theorem~\ref{thm:nopremature}).

%% ============================================================
\section{Related Work}
\label{sec:related}

\paragraph{Proof Engineering in ITP.}
The challenge of maintaining large formal developments is well-documented.
Gross et~al.~\cite{GrossEtAl2024} address proof-engine scaling in Coq through
performant rewriting primitives. Sarracino et~al.~\cite{CoqSurvey2025} survey
Coq users and identify proof brittleness and workflow safety as persistent
community concerns. Isabelle's session model~\cite{IsabelleSessions}
provides build-level governance through document preparation and dependency
tracking. Our work operates at a different granularity: we govern the
\emph{epistemic status} of claims, not the build process, though both
concerns are complementary.

\paragraph{Certificate-Based Approaches.}
Proof certificates appear in several contexts: Coq's
\texttt{Qed} command verifies a proof term against a type;
SAT/SMT solvers produce checkable resolution traces; and
proof-producing compilers (e.g., CakeML) carry correctness
witnesses through compilation.
Our certificates are coarser-grained: they classify evidence
\emph{type} rather than carrying the evidence itself.
This tradeoff sacrifices completeness for lightweight governance
that can span heterogeneous proof methods (formal,
computational, conditional, heuristic).

\paragraph{AI-Assisted Mathematics.}
Machine learning systems for theorem proving---including
tactic prediction~\cite{BlaauwbroekUrbanGeuvers2020} and
autoformalization---are increasingly powerful but can
produce plausible yet unsound output, a phenomenon studied
under ``hallucination''~\cite{FarrugiaEtAl2024}.
Our kernel is source-agnostic: it enforces that
AI-generated heuristic evidence cannot be silently promoted
to proven status, regardless of how convincing it appears.

%% ============================================================
\section{Artifact and Reproducibility}
\label{sec:artifact}

The Lean~4 formalization is available at
\texttt{[repository URL]}\footnote{To be provided upon publication.}.
The artifact comprises:

\begin{itemize}
\item \texttt{Governance.lean}: definitions and core lock/permission theorems
  (82 lines).
\item \texttt{Lambda.lean}: signature classification and case study I
  (148 lines).
\item \texttt{Soundness.lean}: safety theorems (98 lines).
\end{itemize}

\noindent
\textbf{Build.} Install Lean~4 v4.14.0 via \texttt{elan}, then run
\texttt{lake build}. The toolchain is pinned via
\texttt{lean-toolchain} (content: \texttt{leanprover/lean4:v4.14.0}).
The build is deterministic given this toolchain and completes with
zero errors.

\smallskip\noindent
\textbf{Sorry isolation.}
The single \texttt{sorry} in \texttt{Lambda.lean} marks the
analytic content of the ABC argument, which is conditional by
design. It is not imported or used by \texttt{Soundness.lean};
all five kernel safety theorems (Theorems~1--5) are
\texttt{sorry}-free.

%% ============================================================
\section{Conclusion and Future Work}
\label{sec:conclusion}

We have presented a proof governance kernel that enforces typed certificates
and locks on mathematical proof campaigns, preventing unsound claim promotion.
Five safety invariants are machine-checked in Lean~4 with no additional axioms.

Future directions include:
\begin{itemize}
\item \textbf{Fine-grained conditionality.} Extending the
  certificate hierarchy to distinguish GRH-conditional,
  ABC-conditional, and computational results, with composition
  rules that track maximum conditionality through chains.
\item \textbf{Integration with proof assistants.} Embedding the kernel as a
  Lean~4 meta-programming plugin that automatically classifies tactic-generated
  proofs.
\item \textbf{Empirical evaluation.} Applying the kernel to a real proof
  campaign (e.g., the Liquid Tensor Experiment or Fermat's Last Theorem
  formalization) and measuring whether it catches misclassified claims.
\end{itemize}

%% ============================================================
\begin{thebibliography}{10}

\bibitem{BlaauwbroekUrbanGeuvers2020}
L.~Blaauwbroek, J.~Urban, H.~Geuvers.
\newblock Tactic learning and proving for the {C}oq proof assistant.
\newblock In: \emph{LPAR-23}, EPiC Series, 2020.

\bibitem{CoqSurvey2025}
J.~Sarracino et~al.
\newblock Lessons for interactive theorem proving researchers from a survey of
  {C}oq users.
\newblock \emph{J.\ Autom.\ Reason.}, 2025.

\bibitem{DarmonGranville1995}
H.~Darmon, A.~Granville.
\newblock On the equations $z^m = F(x,y)$ and $Ax^p + By^q = Cz^r$.
\newblock \emph{Bull.\ London Math.\ Soc.}, 27(6):513--543, 1995.

\bibitem{Faltings1983}
G.~Faltings.
\newblock Endlichkeitss\"atze f\"ur abelsche {V}ariet\"aten \"uber
  {Z}ahlk\"orpern.
\newblock \emph{Invent.\ Math.}, 73(3):349--366, 1983.

\bibitem{GrossEtAl2024}
J.~Gross, A.~Chlipala, et~al.
\newblock Towards a scalable proof engine: A performant prototype rewriting
  primitive for {C}oq.
\newblock \emph{J.\ Autom.\ Reason.}, 2024.

\bibitem{IsabelleSessions}
{Isabelle Team}.
\newblock Isabelle system manual: Theory sessions.
\newblock \url{https://proofcraft.systems/isabelle/dist/library/Doc/System/Sessions.html}, 2025.

\bibitem{Masser1985}
D.~W.~Masser.
\newblock Open problems.
\newblock In: \emph{Proc.\ Sympos.\ Analytic Number Theory}, London, 1985.

\bibitem{Oesterle1988}
J.~Oesterl\'e.
\newblock Nouvelles approches du ``th\'eor\`eme'' de {F}ermat.
\newblock \emph{S\'em.\ Bourbaki}, exp.~694, 1988.

\bibitem{RatcliffeSurvey2025}
D.~Ratcliffe, B.~Grechuk.
\newblock Generalised {F}ermat equation: A survey of solved cases.
\newblock \emph{Exposition.\ Math.}, 43(4), 2025.

\bibitem{FarrugiaEtAl2024}
S.~Farquhar, J.~Kossen, L.~Kuhn, Y.~Gal.
\newblock Detecting hallucinations in large language models using
  semantic entropy.
\newblock \emph{Nature}, 630:625--630, 2024.

\bibitem{DahmenMPIM2025}
S.~Dahmen.
\newblock Generalized {F}ermat equations with three distinct prime exponents.
\newblock MPIM Number Theory Lunch Seminar, Apr.~2, 2025.

\end{thebibliography}

\end{document}
